problem:  
Let \( M^n \) be a compact, simply connected Riemannian manifold with \( K>0 \).  
Define the *global pinching constant*  
\[
\delta(M) = \frac{\min_{p,\sigma} K_p(\sigma)}{\max_{p,\sigma} K_p(\sigma)},
\]
where \(K_p(\sigma)\) denotes the sectional curvature of the 2-plane \(\sigma \subset T_pM\).  

The **Differentiable Sphere Theorem** of Brendle–Schoen (2009) implies that if \(M\) is *strictly 1/4-pinched in the pointwise sense*, then \(M\) is diffeomorphic to a spherical space form (and to \(S^n\) if simply connected).  

**Problem:**  
Does there exist a constant \(\varepsilon_n>0\), depending only on the dimension \(n\), such that if  
\[
\delta(M) > \frac{1}{4} - \varepsilon_n,
\]  
then \(M\) must be diffeomorphic to the standard sphere \(S^n\)?  

A **related formulation** asks: given a sequence \((M_i^n,g_i)\) of compact simply connected Riemannian manifolds with \(K_i>0\) and \(\delta(M_i) \to \frac{1}{4}\), must every such sequence with \(M_i \not\cong S^n\) (as smooth manifolds) exhibit geometric degeneration—such as unbounded curvature, vanishing injectivity radius, or development of singular spaces in the Cheeger–Gromov limit?

Why is it a "good" problem:  
This problem sharpens our understanding of the **rigidity frontier** in Riemannian geometry. The Brendle–Schoen theorem provides a sharp \(1/4\)-pinching threshold for smooth spherical topology, yet it remains unknown whether even a slight relaxation below this value preserves diffeomorphism to the standard sphere. Resolving this would clarify the **stability of geometric rigidity under curvature perturbations**, testing how delicately curvature controls smooth structure.  

The problem bridges several major areas: positive curvature geometry, geometric analysis via the Ricci flow, and differential topology of exotic spheres. It asks whether nearly quarter-pinched metrics can occur on exotic manifolds or whether such attempts must collapse geometrically. The statement is concise, natural, and tied to a landmark theorem, yet its resolution would require profound advances across multiple fields—a quintessential example of a “good,” research-level mathematics problem.